% chapter6.tex -- Approaching Faces Differently – Pruning Strategy
% Based on the author's UGThesis.docx.

\chapter{Approaching Faces Differently – Pruning Strategy}
\label{ch:pruning_strategy}

In this chapter we develop a new approach to handle biconnected plane graphs by processing faces one after another and pruning branches that inevitably lead to multi-edges. 
We first explain why the original method of processing all faces simultaneously fails, then introduce a safe pruning strategy based on the observation that certain chords persist throughout an entire subtree of the genealogical tree.

\section{Why We Cannot Process All Faces at Once}
\label{sec:why_sequential}

As argued in Chapter~\ref{ch:paper_limitation}, the multi-edge problem arises when two faces share a separation pair and are triangulated independently. 
If we could guarantee that the same chord is never added in two different faces, the original algorithm would work. 
However, because we cannot a priori know which combinations of face triangulations are compatible, we must explore combinations systematically.

One naive approach would be to generate all possible triangulations for each face independently (using the cycle algorithm) and then combine them, discarding those that produce multi-edges. 
This would lead to an exponential blow-up because the total number of combinations is the product of the numbers of triangulations per face. 
Even worse, most combinations might be invalid, wasting time.

Therefore we need a method that explores the space of face triangulations in a coordinated way, building partial triangulations incrementally and backtracking when a conflict is detected.

\section{Hashing Approach and Its Problems}
\label{sec:hashing_problems}

A first idea is to maintain a global set \(S\) of all chords that have already been added during the construction of a partial triangulation. 
We could implement \(S\) as a hash set, which supports insertion, deletion, and membership queries in expected \(O(1)\) time.

When we triangulate a new face, we could check each potential new chord against \(S\) before adding it. 
If a chord already exists, we could skip it or backtrack. 
However, this approach faces two major obstacles:

\begin{enumerate}
    \item \textbf{Invalid triangulations would still be generated.} 
    Even if we check the new chord, the process of flipping chords inside a face may create intermediate triangulations that are invalid (because they contain a chord already in \(S\)). 
    Generating them wastes time, and discarding them after generation does not improve the asymptotic complexity.
    
    \item \textbf{Checking validity takes more than \(O(1)\) time per flip.} 
    If we allow invalid triangulations to be generated, we must check not only the new chord but also whether any previously existing invalid chord has been removed. 
    In general, this would require multiple hash lookups per flip, breaking the constant-time guarantee.
\end{enumerate}

Thus a simple hash set without pruning is insufficient for achieving \(O(1)\) amortised time per valid triangulation.

\section{The Invalid Triangulation Problem}
\label{sec:invalid_problem}

Suppose we generate a triangulation \(T\) of the current face that contains a chord \(e\) which already belongs to \(S\) (i.e., a chord from a previously triangulated face). 
Since we are only allowed to add chords that do not already exist, \(T\) is invalid. 
We could simply skip outputting \(T\) and continue. 
However, the genealogical tree of the current face would still include \(T\) and its descendants. 
If we do not prune the branch, we would generate many descendants, all of which might also be invalid because \(e\) persists.

The key question is: can we safely prune a branch at an invalid node without losing any valid triangulation that might appear deeper in that branch? 
If the answer is yes, we can avoid wasting time on entire subtrees.

\section{Key Observation from the Cycle Algorithm (Persistent Chords)}
\label{sec:persistent_chords}

Recall from Chapter~\ref{ch:cycle_algorithm} that in the genealogical tree of a cycle with root vertex \(v_1\), we only flip chords that are incident to \(v_1\) (the generating chords). 
Chords that do not have \(v_1\) as an endpoint are never flipped; once they appear, they remain in all descendants of that node.

\begin{lemma}[Persistence of non-root chords]
    \label{lem:persistence}
    Let \(T\) be a triangulation of a cycle with root vertex \(v_1\), and let \(e = (v_i, v_j)\) be a chord of \(T\) such that \(i \neq 1\) and \(j \neq 1\). 
    Then for every descendant \(T'\) of \(T\) in the genealogical tree (i.e., every triangulation reachable from \(T\) by a sequence of flips of generating chords), the chord \(e\) belongs to \(T'\).
\end{lemma}

\begin{proof}
    The only operation that removes a chord is flipping it. 
    In the genealogical tree, we only flip generating chords, i.e., chords of the form \((v_1, v_k)\). 
    Since \(e\) does not have \(v_1\) as an endpoint, it is never flipped. 
    Therefore, once \(e\) appears in a triangulation, it persists in all descendants.
\end{proof}

This property is crucial for pruning: if a non-root chord \(e\) conflicts with the global set \(S\) (i.e., it creates a multi-edge), then every descendant triangulation of the current node also contains \(e\) and is therefore invalid. 
Thus we can safely prune the entire subtree rooted at \(T\).

\section{Safe Pruning Condition}
\label{sec:safe_pruning}

Based on Lemma~\ref{lem:persistence}, we formulate a safe pruning condition.

\begin{definition}[Prunable node]
    \label{def:prunable}
    Let \(T\) be a triangulation of the current face \(F_i\) during a depth-first exploration of the global recursion tree. 
    Let \(S\) be the global chord set containing all chords from previously processed faces \(F_1, \dots, F_{i-1}\) and the partial triangulation of face \(F_i\) up to the current node. 
    Suppose that \(T\) contains a chord \(e\) that does not have the root vertex \(v_r\) as an endpoint, and that \(e\) already belongs to \(S\). 
    Then \(T\) is called \textbf{prunable}, and the entire subtree rooted at \(T\) in the local genealogical tree can be skipped without missing any valid complete triangulation.
\end{definition}

\begin{proof}[Justification]
    By Lemma~\ref{lem:persistence}, every descendant \(T'\) of \(T\) also contains \(e\). 
    Since \(e \in S\) and \(e\) is already present in the global chord set (from a previous face), adding \(T'\) would create a multi-edge. 
    Hence all descendants are invalid. 
    No valid complete triangulation can be reached through this branch, so we may prune it.
\end{proof}

\begin{figure}[htbp]
    \centering
    \begin{tikzpicture}[level distance=1.5cm, sibling distance=2cm, scale=0.9, transform shape]
        \tikzset{every node/.style={rectangle, draw, rounded corners, minimum width=1.2cm, minimum height=0.6cm, align=center, font=\scriptsize}}
        \node[fill=blue!20] {Root $T_r$}
            child { node[fill=green!15] {$T_1$ (valid)}
                child { node[fill=yellow!15] {$T_{11}$ (valid)} }
                child[red] { node[fill=red!30] {$T_{12}$ (invalid, $e \in S$)} 
                    child { node[fill=gray!20] {$T_{121}$ (pruned)} }
                    child { node[fill=gray!20] {$T_{122}$ (pruned)} }
                }
            }
            child { node[fill=green!15] {$T_2$ (valid)} };
        \draw[red, dashed] (T12.south) to[out=-90,in=90] (T121.north);
        \draw[red, dashed] (T12.south) to[out=-90,in=90] (T122.north);
        \node[red] at (4, -2.5) {Pruned subtree};
    \end{tikzpicture}
    \caption{Safe pruning: node $T_{12}$ contains a chord $e$ already in $S$; its entire subtree is pruned because all descendants would also contain $e$.}
    \label{fig:pruning}
\end{figure}

\section{Requirements for Safe Pruning}
\label{sec:pruning_requirements}

To implement this pruning strategy in a full algorithm for biconnected plane graphs, we must solve the following problems:

\begin{enumerate}
    \item Process faces one by one in a fixed order so that the global set \(S\) is well-defined.
    \item Ensure that chords from previous faces are never removed (they are permanent). 
    \item During the triangulation of a face, maintain a local genealogical tree where flips only affect chords incident to the root vertex.
    \item Check, before flipping a generating chord, whether the new chord (the opposite pair) already exists in \(S\). 
    If it does, skip that flip (prune the branch).
    \item When a chord that does not have the root vertex as an endpoint is added (as a result of a flip), we must check whether it conflicts with \(S\). 
    If it does, we prune the entire subtree from that node onward.
    \item Guarantee that no valid triangulation is missed by this pruning.
    \item Ensure that all operations (insertion, deletion, membership) on \(S\) take \(O(1)\) expected time, and that the number of pruned nodes does not dominate the total work.
\end{enumerate}

\section{Outline of the Solution}
\label{sec:pruning_outline}

In the following chapters we address these requirements:

\begin{itemize}
    \item \textbf{Chapter~\ref{ch:multi_layer_dfs}} presents the multi-layer DFS that processes faces sequentially and maintains the global chord set \(S\) with push/pop operations for backtracking.
    \item \textbf{Chapter~\ref{ch:generating_chord_conflicts}} discusses the problem of generating chords themselves causing multi-edges and introduces the concept of a \emph{safe root vertex}.
    \item \textbf{Chapter~\ref{ch:safe_root_proof}} proves that a safe root always exists and gives an \(O(n)\) algorithm to find one.
    \item \textbf{Chapter~\ref{ch:lower_bound}} establishes that each face yields at least \(\Omega(n)\) valid triangulations, ensuring amortised constant time.
    \item \textbf{Chapter~\ref{ch:vgs}} introduces the valid generating set (VGS) and constant-time updates.
\end{itemize}

\section{Summary}
\label{sec:pruning_summary}

We have identified the persistence property of non-root chords as the key to safe pruning. 
If a triangulation of a face contains a chord that conflicts with previous faces, that chord will remain in all descendants, making the entire subtree invalid. 
Therefore we can prune such branches without losing any valid triangulation. 
This observation allows us to process faces sequentially while maintaining optimal complexity, provided we can also handle conflicts that involve generating chords themselves (which are root-incident and thus can be flipped away). 
The next chapter describes the multi-layer DFS that implements this pruning strategy.